\section{Resolution towers and the regime taxonomy}

This section records claim nodes and dependency edges.

\begin{proposition}[The resolution tower]\label{thm:resolution-tower}\lean{Obsmetrologia.Closure.ResolutionTower.tower}\notready\uses{thm:sampling-fibration, thm:pixelization-closure, thm:reflexive-tower}
\emph{(Open frontier --- stated, not formalized.)} A refinement tower of sampling
fibrations $\texttt{base}_{n} : \texttt{Domain} \to \texttt{GridPoint}_{n}$, each stage
strictly finer, in the pattern of \cref{thm:reflexive-tower}, with the stagewise
pixelization closures \cref{thm:pixelization-closure} assembled into one tower closure.
At any \emph{single} grid the closure converges in one application --- trivially strict,
$\kappa = 1$; that degeneracy is recorded here honestly, and it is why the tower, not
the single grid, is the meaningful object. Naturality across stages
($\texttt{pixelize}\,F_{n} \circ \texttt{refine} = \texttt{coarsen} \circ
\texttt{pixelize}\,F_{n+1}$, in the commuting-square idiom) is part of this node; per
the dervish precedent, no strong-form functor node is attempted until the leaves are
green.
\end{proposition}

\begin{theorem}[Band-limitedness as regime membership]\label{thm:bandlimited-regime}\lean{Obsmetrologia.Closure.ResolutionTower.bandlimited_regime}\notready\uses{thm:resolution-tower, thm:four-regime-dichotomy, thm:strict-finite, thm:kappa-stratification-interface}
\emph{(Open frontier --- stated, not formalized.)} The tower closure's stabilization
stage read through \cref{thm:four-regime-dichotomy}: stabilization at a finite stage is
the \texttt{strict} regime (piecewise-constant imagery at some finite resolution);
stabilization only at $\omega$ is the \texttt{scott} regime (classical band-limitedness,
recovered only as sampling density grows without bound --- Nyquist--Shannon's own
framing); the transfinite and aperiodic regimes are mathematically available and
physically unclaimed, recorded as such. Prior art (Haralick 1989; Heijmans 1992; Unser--Aldroubi)
is uniformly finite-lattice or Hilbert-space and never sees the ordinal-height question;
this node is the chapter's claim to new territory.
\end{theorem}
